% Nouveau EGA V (cote 25): archive cote 25, site page 113 (PDF page 114).
% typescript. Draft compared: none.
% AI-transcribed by Claude Opus 5.5 (claude-opus-5-5, adaptive thinking, effort high), Message Batches API, batch ega25b (msgbatch_01A7n5YuXubUPNK1zsepi8RG), 2026-11-13; checked against both supplied scan parts; see CHECK_DECISIONS.tsv.
\scanpage{25/113}
\mspage{\hand{16}}

\xout{Par symétrie,} Lorsque Z est fermé, alors X-Z est le plus grand ouvert
tel que la famille de diviseurs de $\overset{\hand{\unsure{r}}}{X}-Z$ paramétrée par T, induite par la
famille donnée (dans un sens évident) est sans point fixe.

Si la famille D est sans point fixe, et si T est \inserted[typed]{plat} loc de prés finie
sur S à fibres $(S_1)$\inserted[typed]{et géom irréd}, alors D est aussi un diviseur \deleted{de Cartier}
relativement à X (pour $pr_1 : X\times_S T \hand{\longrightarrow} X$): en effet, D est défini localemen
% NOTE: the right edge of the page is cut; the end of this line and of several following lines is not visible
\inserted[typed]{au point $z \hand{\in}$ supp D} par une équation \xout{f=0} $\hand{\phi} = 0$, et l'équation induite sur la fibre
$T_{s\ k(x)}$ du point $x\hand{\in} X$ (au dessus de $s\hand{\in} S$) est non \xout{identiquement nulle}
nilpo\xout{n}tente en z (car autrement $pr_1^{-1}(x)\ D = V(\hand{\phi}_x)$ \xout{serait} \inserted[typed]{contiendrait} ensemblist\xout{\unsure{e}}
ment un voisinage de z dans $T_{\xout{s}\ k(x)}$ , donc contiendrait $T_{s\ k(x)}$ , i.e.
x serait point fixe, ce qui n'est pas);or comme $T_{s\ k(x)}$ est irréd et $(S_1)$
il s'ensuit que $\hand{\phi}_x$ est $\underline{O}_{T_{\xout{s}\ k(\unsure{x})}}$-régulière en z . On obtient ainsi
\xout{un morphisme} une famille de div de T/S paramétrée par X, i.e. un morph
\[ X \hand{\longrightarrow} \mathrm{Div}_{T/S} \quad \hand{,} \]
\over{D}{D}ans le cas général où la famille de \unsure{diviseurs} de X peut avoir des points
fixes, on obtient une famille de diviseurs de T/S paramétrée par X-Z,
i;e. un morphisme
\[ X - Z \hand{\longrightarrow} \mathrm{Div}_{T/S} \ , \]
en remplaçant dans la définition précédente X pa\unsure{r|m} X-T . D'ailleurs la
\xout{si D est sans points fixes,}

\xout{Notons en passant que la famille de diviseurs (notée ${}^{t}D$ )}
\xout{X-Z}
\xout{de T/S paramétrée par X }\deleted{\xout{est év}idemment sans points fixes, et on}
\xout{obtient ainsi une correspondance biunivo}\deleted{que entre familles de diviseurs}
\xout{sans points fixes} démonstration ci-dessus montre que X-Z est exactement
le plus grand ouvert \inserted[typed]{U} de X tel que $D\hand{|}(U\times_S T)$ soit un di\xout{vi}seur \deleted{de Cartier}
relativement à U, i.e. \inserted[typed]{tel que son symétrique ${}^{t}D$} \unsure{``soit''} une famille de diviseurs de T/S paramétrée
par U/S .

\mspage{\hand{113}}
